\chapter*{คำนำ}
\addcontentsline{toc}{chapter}{คำนำ}

ในศตวรรษที่ 21 โลกกำลังเผชิญกับความท้าทายครั้งใหญ่ที่มีความเชื่อมโยงกันอย่างแยกไม่ออก นั่นคือการสร้างความมั่นคงทางอาหาร (Food Security) ควบคู่ไปกับการรักษาสมดุลทางนิเวศวิทยาและการฟื้นฟูสิ่งแวดล้อม (Environmental Sustainability) ภาคการเกษตรดั้งเดิมที่พึ่งพาปุ๋ยเคมีสังเคราะห์ ยาปราบศัตรูพืช และพลังงานฟอสซิลอย่างเข้มข้น ไม่เพียงแต่ส่งผลให้เกิดการเสื่อมโทรมของทรัพยากรดินและน้ำ หากยังเป็นแหล่งปลดปล่อยก๊าซเรือนกระจกที่สำคัญ ยิ่งไปกว่านั้น ปัญหามลพิษทางอากาศ ขยะอุตสาหกรรม และการปนเปื้อนของสารเคมีอันตรายที่ย่อยสลายยากในสิ่งแวดล้อม ได้ทวีความรุนแรงจนส่งผลกระทบต่อสุขภาพของมนุษย์และสิ่งมีชีวิตในระบบนิเวศอย่างไม่อาจหลีกเลี่ยงได้

ท่ามกลางวิกฤตการณ์เหล่านี้ การบูรณาการองค์ความรู้ข้ามศาสตร์ (Interdisciplinary Integration) ได้กลายเป็นกุญแจสำคัญในการแสวงหาทางออกใหม่ที่ยั่งยืน \textbf{เทคโนโลยีพลาสมา (Plasma Technology)} ซึ่งเดิมทีเป็นศาสตร์ทางฟิสิกส์ขั้นสูงที่ถูกจำกัดอยู่ในวงการดาราศาสตร์ พลังงานนิวเคลียร์ฟิวชัน อุตสาหกรรมเซมิคอนดักเตอร์ และวิศวกรรมวัสดุศาสตร์ ได้ก้าวข้ามขีดจำกัดเดิมผ่านการพัฒนาเทคโนโลยีพลาสมาความดันบรรยากาศอุณหภูมิต่ำ หรือพลาสมาเย็น (Cold Atmospheric Plasma หรือ CAP) ที่สามารถทำงานได้ในสภาวะแวดล้อมปกติโดยไม่ต้องใช้ระบบสุญญากาศที่ซับซ้อนและสิ้นเปลืองพลังงาน

ความก้าวหน้าดังกล่าวได้เปิดพรมแดนใหม่ที่เรียกว่า \textbf{พลาสมาเพื่อการเกษตร (Plasma Agriculture)} และ \textbf{พลาสมาเพื่อการจัดการสิ่งแวดล้อม (Environmental Plasma Technology)} โดยอาศัยกลไกอันเป็นเอกลักษณ์ของพลาสมาในการผลิตอนุภาคที่มีความว่องไวทางเคมีสูง โดยเฉพาะอนุมูลอิสระของออกซิเจนและไนโตรเจน (Reactive Oxygen and Nitrogen Species หรือ RONS) สนามแม่เหล็กไฟฟ้า และรังสียูวี โดยใช้เพียงอากาศ น้ำ และพลังงานไฟฟ้าเป็นวัตถุดิบตั้งต้น เทคโนโลยีนี้จึงเปรียบเสมือนโรงงานเคมีสีเขียวขนาดจิ๋วที่สามารถกระตุ้นความงอกของเมล็ดพันธุ์ ตรึงไนโตรเจนจากอากาศมาสร้างเป็นปุ๋ยน้ำ (Plasma-Activated Water หรือ PAW) ยับยั้งเชื้อก่อโรคพืชและจุลินทรีย์ปนเปื้อนในอาหารหลังการเก็บเกี่ยว ตลอดจนย่อยสลายสารพิษตกค้างและมลพิษในน้ำ ดิน และอากาศได้อย่างมีประสิทธิภาพโดยไม่ทิ้งสารเคมีตกค้าง

หนังสือ \textbf{พลาสมาฟิสิกส์และปัญญาประดิษฐ์ เพื่อการเกษตรและสิ่งแวดล้อม: จากฟิสิกส์การคายประจุและการเรียนรู้เชิงลึกสู่นวัตกรรมยั่งยืน} เล่มนี้ เรียบเรียงขึ้นโดยมีวัตถุประสงค์เพื่อเป็นสะพานเชื่อมองค์ความรู้ระหว่างนักฟิสิกส์ วิศวกร นักวิทยาศาสตร์การเกษตร นักจุลชีววิทยา และนักวิทยาศาสตร์สิ่งแวดล้อม เนื้อหาภายในเล่มได้รับการออกแบบให้มีความต่อเนื่องอย่างเป็นระบบ เริ่มตั้งแต่การปูพื้นฐานทางฟิสิกส์ของสถานะพลาสมาและเคมีของการคายประจุในก๊าซด้วยภาษาที่เข้าใจง่าย ไปจนถึงการเจาะลึกกลไกระดับเซลล์และโมเลกุลในการประยุกต์ใช้งานทางการเกษตรตั้งแต่ต้นน้ำถึงปลายน้ำ การบำบัดมลพิษสิ่งแวดล้อมครบวงจร และการออกแบบขยายขนาดระบบสู่การใช้งานจริงในระดับอุตสาหกรรม

ผู้เขียนหวังเป็นอย่างยิ่งว่า หนังสือเล่มนี้จะเป็นประโยชน์ต่อนักศึกษาระดับปริญญาตรีและบัณฑิตศึกษา คณาจารย์ นักวิจัย ตลอดจนผู้ประกอบการด้านเทคโนโลยีเกษตรสมัยใหม่ (AgriTech) และเทคโนโลยีสะอาด (CleanTech) ในการใช้เป็นตำราเรียน เอกสารอ้างอิง และแรงบันดาลใจในการสร้างสรรค์งานวิจัยและนวัตกรรมพลาสมา เพื่อร่วมกันขับเคลื่อนภาคการเกษตรและการจัดการสิ่งแวดล้อมของประเทศไทยให้ก้าวสู่ความยั่งยืนอย่างแท้จริง

\vspace{1.5cm}
\begin{flushright}
    \textbf{ผู้ช่วยศาสตราจารย์ ดร.ชีวะ ทัศนา}\\
    สาขาวิชาปัญญาประดิษฐ์เพื่อเกษตรดิจิทัลและสิ่งแวดล้อม\\
    คณะวิทยาศาสตร์และเทคโนโลยี มหาวิทยาลัยราชภัฏรำไพพรรณี\\
    พ.ศ. 2569
\end{flushright}

\clearpage
